\rsection{Wahrnehmung}
\slide{Begrifflichkeiten \& Grundlagen}{
    \begin{itemize}
        \item \b{Reizschwelle} (Threshold): Die schwächste Reizintensität, bei der ein Reiz gerade noch wahrgenommen werden kann.
        \item \b{Bottom-Up-Prozesse:} Reizgesteuerte Verarbeitung; Informationen werden aus dem physikalischen Stimulus aufgenommen und verarbeitet.
        \item \b{Top-Down-Prozesse:} Wissensgesteuerte Verarbeitung; Vorwissen, Erfahrung und Kontext beeinflussen die Wahrnehmung        
        \item \b{Interaktion der Prozesse:} Wahrnehmung ist meist eine Mischung aus Bottom-Up (Daten) und Top-Down (Konzepten).
    \end{itemize}
}
\slide{Gestaltpsychologie}{
    \begin{itemize}
        \item \b{Gestaltprinzipien:} Regeln zur Gruppierung von Reizen, z. B. Nähe, Ähnlichkeit, Geschlossenheit, gute Fortsetzung und gemeinsame Bewegung.
        \item \b{Kontexteffekte:} Reizeigenschaften werden unbewusst mit Kontext und Vorwissen verrechnet (z. B. Helligkeits-, Größen- und Formkonstanz).
        \item \b{Figur-Grund-Differenzierung:} Fähigkeit, zwischen vordergründiger Figur und Hintergrund zu unterscheiden.
    \end{itemize}
}
\slide{Objekterkennung (Theorien)}{
    \begin{itemize}
        \item \b{Multi-Level-Theory (Marr):} Objekterkennung als Informationsverarbeitung auf drei Ebenen (Computational, Algorithmisch, Implementierung).
        \begin{itemize}
            \item Prozess-Stufen: Primal Sketch (Kanten/Konturen) $\rightarrow$ 2,5-D Sketch (Tiefe/Orientierung, egozentrisch) $\rightarrow$ 3D-Modell (objektzentriert, allozentrisch).
        \end{itemize}
        \item \b{Komponententheorie (Biederman):} Objekte werden durch die Identifikation und Zusammensetzung elementarer geometrischer Teilobjekte ("Geone") erkannt.
    \end{itemize}
}
\slide{Modularität der Wahrnehmung (Fodor)}{
    \begin{itemize}
        \item \b{Modularitätshypothese:} Kognitive Prozesse (wie Wahrnehmung) finden in unabhängigen, spezialisierten Modulen statt.
        \item \b{Informational Encapsulation:} Module sind informationell abgekapselt; sie haben keinen Zugriff auf Hintergrundwissen oder andere Module.
        \item Weitere Eigenschaften: Automatisch, schnell, domänen-spezifisch, neuronale Architektur.
        \item Pro Modularität: Optische Täuschungen bleiben bestehen, auch wenn man weiß, dass es Täuschungen sind (Wissen korrigiert Wahrnehmung nicht).
        \item Contra Modularität: Kontexteffekte zeigen, dass Wissen und Erfahrung die Wahrnehmung beeinflussen (Top-Down-Einflüsse).
    \end{itemize}
}